% ============================================================
%  Openbook / Formelsammlung Template — strikt nach my_dhbw
%  Stil: 2-spaltig, kompakt, scriptsize Tabellen, dünne Regeln.
%  Renderer: lualatex (zweimal, wg. Unicode-Inhalt)
% ============================================================
\documentclass[10pt, a4paper]{article}

\usepackage[ngerman]{babel}
\usepackage{fontspec}
\setmainfont[
  Path=/usr/share/fonts/TTF/,
  Extension=.ttf,
  BoldFont=DejaVuSans-Bold,
  ItalicFont=DejaVuSans-Oblique,
  BoldItalicFont=DejaVuSans-BoldOblique
]{DejaVuSans}
\setsansfont[
  Path=/usr/share/fonts/TTF/,
  Extension=.ttf,
  BoldFont=DejaVuSans-Bold,
  ItalicFont=DejaVuSans-Oblique,
  BoldItalicFont=DejaVuSans-BoldOblique
]{DejaVuSans}
\setmonofont[Path=/usr/share/fonts/TTF/,Extension=.ttf]{DejaVuSansMono}

\usepackage[top=0.8cm, bottom=0.8cm, left=0.8cm, right=0.8cm, footskip=0.4cm]{geometry}
\usepackage{amsmath, amssymb}
\usepackage{multicol}
\usepackage{xcolor}
\usepackage{titlesec}
\usepackage{enumitem}
\usepackage{fancyhdr}
\usepackage{lastpage}
\usepackage{microtype}
\usepackage{etoolbox}

\usepackage{booktabs}
\usepackage{array}
\usepackage{tabularx}
% ltablex/longtable bewusst NICHT geladen — kollidiert mit multicols.
\usepackage{colortbl}
\usepackage[most]{tcolorbox}
\usepackage{listings}
\usepackage[normalem]{ulem}
\usepackage{adjustbox}
\usepackage[colorlinks=false]{hyperref}

\renewcommand{\familydefault}{\sfdefault}

% ── Theme (my_dhbw accent) ────────────────────────────────────
\definecolor{accent}{RGB}{0, 60, 113}
\definecolor{primaryblue}{RGB}{0, 60, 113}
\definecolor{codebg}{RGB}{245, 245, 245}
\definecolor{figurebg}{RGB}{232, 241, 255}
\definecolor{tablehintbg}{RGB}{240, 255, 240}
\definecolor{solutionbg}{RGB}{242, 255, 242}
\definecolor{rulegray}{RGB}{180, 180, 180}
\definecolor{tableheadbg}{RGB}{0, 60, 113}

% ── Header/Footer — wie my_dhbw formelsammlung.tex ────────────
\pagestyle{fancy}
\fancyhf{}
\renewcommand{\headrulewidth}{0pt}
\renewcommand{\footrulewidth}{0.3pt}
\fancyfoot[L]{\tiny\color{accent}\nichttitle}
\fancyfoot[R]{\tiny\color{accent}\thepage\,/\,\pageref{LastPage}}

% ── Typografie — wie formelsammlung.tex ───────────────────────
\titleformat{\section}
    {\color{accent}\footnotesize\bfseries}{}{0pt}{}
    [\vspace{-2pt}\color{accent!50}\rule{\linewidth}{0.4pt}]
\titleformat{\subsection}
    {\scriptsize\bfseries\color{accent!85!black}}{}{0pt}{}{}
\titleformat{\subsubsection}
    {\scriptsize\itshape\color{accent!60!black}}{}{0pt}{}{}
\titlespacing*{\section}{0pt}{4pt}{1pt}
\titlespacing*{\subsection}{0pt}{3pt}{0pt}
\titlespacing*{\subsubsection}{0pt}{2pt}{0pt}

\setlength{\parindent}{0pt}
\setlength{\parskip}{1pt}
\renewcommand{\arraystretch}{1.0}
\setlist{noitemsep, topsep=0pt, parsep=0pt, leftmargin=1.1em}

\setlength{\columnsep}{6mm}
\setlength{\columnseprule}{0.2pt}

% Globaler scriptsize-Body. Tabellen zusätzlich auf scriptsize zwingen.
\AtBeginEnvironment{tabularx}{\scriptsize}
\AtBeginEnvironment{tabular}{\scriptsize}

% ── Boxen (kompakt) ───────────────────────────────────────────
\newtcolorbox{figurebox}[1]{
  colback=figurebg, colframe=accent!50,
  title={\scriptsize Abbildung: #1}, fonttitle=\scriptsize\bfseries,
  fontupper=\scriptsize, sharp corners,
  boxsep=2pt, left=3pt, right=3pt, top=2pt, bottom=2pt,
  before skip=3pt, after skip=3pt
}
\newtcolorbox{tablehintbox}[1]{
  colback=tablehintbg, colframe=green!50!black!50,
  title={\scriptsize Tabelle: #1}, fonttitle=\scriptsize\bfseries,
  fontupper=\scriptsize, sharp corners,
  boxsep=2pt, left=3pt, right=3pt, top=2pt, bottom=2pt,
  before skip=3pt, after skip=3pt
}
\newtcolorbox{solutionbox}[1]{
  colback=solutionbg, colframe=green!60!black!60,
  title={\scriptsize #1}, fonttitle=\scriptsize\bfseries,
  fontupper=\scriptsize, sharp corners,
  boxsep=2pt, left=3pt, right=3pt, top=2pt, bottom=2pt,
  before skip=3pt, after skip=3pt
}

\lstset{
  basicstyle=\ttfamily\tiny,
  backgroundcolor=\color{codebg},
  breaklines=true,
  frame=none,
  xleftmargin=2pt, xrightmargin=2pt,
  aboveskip=2pt, belowskip=2pt,
  keepspaces=true
}

\providecommand{\nichttitle}{Formelsammlung}
\providecommand{\nichtsubtitle}{}

\renewcommand{\nichttitle}{Digitale Betriebswirtschaftslehre}
\renewcommand{\nichtsubtitle}{Formelsammlung}


\begin{document}

\begin{center}
{\small\bfseries\color{accent}\nichttitle}
\end{center}
\vspace{-6pt}

\scriptsize
\begin{multicols}{2}

% ── BODY ──────────────────────────────────────────────────────

\section{Digitale Betriebswirtschaftslehre — Open-Book-Spickzettel}


\subsection{1. Grundlagen}

\textbf{Wirtschaften}: Knappe Güter geplant einsetzen → Bedürfnisbefriedigung möglichst vorteilhaft.

\textbf{Wertschöpfung} = Output − Input | Modell: \textbf{Porters Wertschöpfungskette}

\begin{itemize}
  \item Input (Produktionsfaktoren) → Transformation → Output
\end{itemize}


\subsection{2. Ansätze der BWL}


{\small
\begin{adjustbox}{max width=\linewidth}
{\scriptsize
\begin{tabular}{lllll}
\toprule
\textbf{Ansatz} & \textbf{Vertreter} & \textbf{Blick} & \textbf{Kernfrage} & \textbf{Zielgröße} \\
\midrule
Produktionsorientiert & Gutenberg & Produktion & Wie effizient produzieren? & Produktivität (Output/Input) \\
Entscheidungsorientiert & Heinen & Entscheidungen & Wie gute Entscheidungen? & Qualität von Entscheidungen \\
Systemorientiert & Ulrich & Gesamtsystem & Wie hängt alles zusammen? & Stabilität, Überlebensfähigkeit \\
Institutionenökonomisch & Coase & Regeln/Verträge & Wie effizient koordinieren? & Min. Transaktionskosten \\
\bottomrule
\end{tabular}
}
\end{adjustbox}
}


\begin{itemize}
  \item \textbf{System}: offen, Elemente + Umwelt + Rückkopplungen
  \item \textbf{Institutionenökonomisch}: Transaktionskosten, Informationsasymmetrien, Principal-Agent
\end{itemize}


\subsection{3. Produktionsfaktoren (Gutenberg)}

\textbf{Arbeit} | \textbf{Betriebsmittel (Kapital)} | \textbf{Werkstoffe}

\textbf{Daten} (modern): Sammlung → Verarbeitung → Entscheidung → Wertschöpfung



\subsection{4. Principal-Agent}

\begin{itemize}
  \item \textbf{Principal} = Eigentümer (delegiert) | \textbf{Agent} = Manager (entscheidet)
  \item \textbf{Informationsasymmetrie}: Principal weiß weniger
  \item \textbf{Lösung}: leistungsabh. Vergütung, KPIs, Monitoring
\end{itemize}


\subsection{5. Shareholder vs. Stakeholder}


{\small
\begin{adjustbox}{max width=\linewidth}
{\scriptsize
\begin{tabular}{lll}
\toprule
\textbf{Merkmal} & \textbf{Stakeholder} & \textbf{Shareholder} \\
\midrule
Fokus & alle Anspruchsgruppen & Eigentümer \\
Ziel & Interessen aller & Max. Unternehmenswert \\
Maßstab & Max. Differenz Nutzen − Kosten aller & Max. zukünftige Zahlungen an Eigentümer \\
Beurteilung & nicht operational, pluralistisch & operational, monistisch \\
\bottomrule
\end{tabular}
}
\end{adjustbox}
}



{\small
\begin{adjustbox}{max width=\linewidth}
{\scriptsize
\begin{tabular}{lll}
\toprule
\textbf{Stakeholder} & \textbf{Typ} & \textbf{Anspruch} \\
\midrule
Eigenkapitalgeber & intern & Kapitalmehrung \\
Arbeitnehmer & intern & Entlohnung, Sicherheit \\
Management & intern & Gehalt, Einfluss, Status \\
Fremdkapitalgeber & extern & Tilgung + Verzinsung \\
Kunden & extern & Qualität, Wettbewerbsfähigkeit \\
Lieferanten & extern & zuverlässige Zahlung \\
Staat/Gesellschaft & extern & Steuern, Recht, Umweltschutz \\
\bottomrule
\end{tabular}
}
\end{adjustbox}
}



\subsection{6. Unternehmensziele}

\textbf{Funktionen}: Orientierung, Entscheidungskriterium, Steuerung, Koordination, Legitimation, Kontrolle.



{\small
\begin{adjustbox}{max width=\linewidth}
{\scriptsize
\begin{tabular}{lll}
\toprule
\textbf{} & \textbf{Sachziele} & \textbf{Formalziele} \\
\midrule
Bezug & Konkretes Handeln & Übergeordnete Erfolgsziele \\
Beispiele & Produkte, Qualität, Marktposition & Gewinn, Rentabilität, Liquidität \\
\bottomrule
\end{tabular}
}
\end{adjustbox}
}


\textbf{Sachziel-Kategorien}: Leistungs-, Finanz-, Führungs-/Organisations-, soziale/ökologische Ziele.


\textbf{Zeitbezug}: kurzfristig ≤ 1 J | mittelfristig 1–5 J | langfristig > 5 J

\textbf{Messbarkeit}: kardinal (€) | ordinal (gut/schlecht) | nominal (ja/nein)

\textbf{Rangordnung}: Ober- → Zwischen- → Unterziele (Hierarchie)



\subsubsection{Zielbeziehungen}


{\small
\begin{adjustbox}{max width=\linewidth}
{\scriptsize
\begin{tabular}{ll}
\toprule
\textbf{Typ} & \textbf{Bedeutung} \\
\midrule
Komplementär & A fördert B \\
Konkurrierend (konfliktär) & A behindert B \\
Indifferent & unabhängig \\
\bottomrule
\end{tabular}
}
\end{adjustbox}
}


\textbf{Beispielpaare}:

\begin{itemize}
  \item Personal. Service ↔ Reduktion Filial-Service → konkurrierend
  \item KI-Automatisierung ↔ Reduktion Ausschuss → komplementär
  \item Flexible Arbeitszeit ↔ Vertrieb-Steigerung → indifferent
  \item Just-in-Time ↔ Lagerkosten ↓ → komplementär
  \item Produktivität (digital) ↔ Investitionen verschoben → konkurrierend
  \item Kapazität stilllegen ↔ Absatzsteigerung bei konst. Preis → konkurrierend
\end{itemize}


\subsubsection{SMART}
\textbf{S}pezifisch | \textbf{M}essbar | \textbf{A}ttraktiv | \textbf{R}ealistisch | \textbf{T}erminiert



\subsection{7. Kennzahlen (Formalziele)}

\begin{itemize}
  \item \textbf{Produktivität} = Output-Menge / Input-Menge (mengenmäßig)
  \item \textbf{Wirtschaftlichkeit} = Ertrag (€) / Aufwand (€) (wertmäßig)
  \item \textbf{Gewinn} = Ertrag − Aufwand
  \item \textbf{Rentabilität} = (Gewinn / Ø eingesetztes Kapital) × 100
\end{itemize}


\subsubsection{Beispiel Schraubenproduktion}
\begin{itemize}
  \item 1.000 Schrauben / 10 kg = \textbf{100 Schrauben/kg}
  \item Aufwand 10·2 = 20 € | Ertrag 1000·0,02 = 20 € → \textbf{W = 1,0}
  \item +10 \%: \textbf{110 Schrauben/kg}, \textbf{W = 1,1}
\end{itemize}

\textbf{Hebel}: Output ↑ | Input ↓ | VK-Preis ↑ | EK-Preis ↓



\subsection{8. Ökonomisches Prinzip}


{\small
\begin{adjustbox}{max width=\linewidth}
{\scriptsize
\begin{tabular}{ll}
\toprule
\textbf{Prinzip} & \textbf{Logik} \\
\midrule
Maximum & gegebener Input → max. Output \\
Minimum & gegebener Output → min. Input \\
Optimum & bestes Verhältnis Output/Input \\
\bottomrule
\end{tabular}
}
\end{adjustbox}
}



\subsection{9. Effizienz vs. Effektivität}


{\small
\begin{adjustbox}{max width=\linewidth}
{\scriptsize
\begin{tabular}{lll}
\toprule
\textbf{} & \textbf{Effizienz} & \textbf{Effektivität} \\
\midrule
Verhältnis & Leistung/Ressourcen & Zielerreichungsgrad \\
Frage & Dinge richtig machen? & Richtige Dinge machen? \\
Begriff & Leistungsfähigkeit & Leistungswirksamkeit \\
\bottomrule
\end{tabular}
}
\end{adjustbox}
}



\subsection{10. Mindestbedingungen}

\begin{itemize}
  \item \textbf{Langfristig (Erfolg)}: Erträge ≥ Aufwendungen
  \item \textbf{Kurzfristig (Liquidität)}: Einzahlungen ≥ Auszahlungen
\end{itemize}


\subsection{11. Markt — Grundlagen}

\textbf{Markt}: Ort, an dem Angebot + Nachfrage aufeinandertreffen.

\begin{itemize}
  \item \textbf{Nachfrage}: Preis ↑ → QD ↓
  \item \textbf{Angebot}: Preis ↑ → QS ↑
  \item \textbf{Grenzen}: Informationsasymmetrien, Markteintrittsbarrieren, Marktmacht
\end{itemize}


\subsection{12. Marktgleichgewicht}

\textbf{Gleichgewicht}: QD = QS → markträumend.

\begin{itemize}
  \item QD > QS → Nachfrageüberschuss → Preis ↑
  \item QS > QD → Angebotsüberschuss → Preis ↓
\end{itemize}


\subsubsection{Beispiel Ladekabel: QD = 100 − 2P | QS = 20 + 2P}
\begin{enumerate}
  \item 100 − 2P = 20 + 2P → 80 = 4P → \textbf{P\textbackslash{}* = 20 €}
  \item \textbf{Q\textbackslash{}* = 60}
\end{enumerate}


\subsubsection{Disequilibrium P = 10 €}
\begin{itemize}
  \item QD = 80 | QS = 40 → Nachfrageüberschuss 40 → Preis steigt bis P\textbackslash{}* = 20
\end{itemize}


\subsubsection{Nachfrageverschiebung QD = 120 − 2P}
\begin{itemize}
  \item 100 = 4P → \textbf{P\textbackslash{}* = 25 €}, \textbf{Q\textbackslash{}* = 70} (Rechtsverschiebung → Preis ↑, Menge ↑)
  \item Ableitung: Kapazität ausbauen, Preis anheben
\end{itemize}

\begin{quote}
Merke: Datenbasierte Nachfrageschätzung (Verkaufszahlen, Suchanfragen) → schnellere Preis-/Mengenanpassung.
\end{quote}


\subsection{13. Marktformen}


{\small
\begin{adjustbox}{max width=\linewidth}
{\scriptsize
\begin{tabular}{llll}
\toprule
\textbf{} & \textbf{viele Anbieter} & \textbf{wenige Anbieter} & \textbf{ein Anbieter} \\
\midrule
\textbf{viele Nachfrager} & Polypol & Oligopol & Monopol \\
\textbf{wenige Nachfrager} & Nachfrageoligopol & zweiseitiges Oligopol & beschränktes Monopol \\
\textbf{ein Nachfrager} & Nachfragemonopol & beschränktes Monopol & zweiseitiges Monopol \\
\bottomrule
\end{tabular}
}
\end{adjustbox}
}



{\small
\begin{adjustbox}{max width=\linewidth}
{\scriptsize
\begin{tabular}{llll}
\toprule
\textbf{Merkmal} & \textbf{Polypol} & \textbf{Oligopol} & \textbf{Monopol} \\
\midrule
Anbieter & viele & wenige & einer \\
Marktmacht & keine & mittel & hoch \\
Wettbewerb & sehr hoch & begrenzt & gering \\
\bottomrule
\end{tabular}
}
\end{adjustbox}
}



\subsection{14. Wettbewerb}

\begin{itemize}
  \item Strategien: \textbf{Kostenführerschaft} | \textbf{Differenzierung} | \textbf{Fokussierung}
  \item \textbf{Markteintrittsbarrieren}: hohe Investitionskosten, Skaleneffekte, Vertriebskanäle, Regulierung → weniger Wettbewerb, mehr Marktmacht
\end{itemize}


\subsection{15. Digitale Märkte \& Plattformen}

\begin{itemize}
  \item \textbf{Plattformmärkte}: vermitteln ≥ 2 Nutzergruppen (Amazon, Uber, Airbnb)
  \item \textbf{Netzwerkeffekte}: mehr Nutzer → höherer Nutzen → Marktkonzentration
  \item \textbf{Lock-in}: Wechselkosten/Gewohnheit binden Kunden
\end{itemize}


\subsection{16. Entscheidung — Arten}


{\small
\begin{adjustbox}{max width=\linewidth}
{\scriptsize
\begin{tabular}{lll}
\toprule
\textbf{Art} & \textbf{Merkmale} & \textbf{Beispiele} \\
\midrule
Konstitutiv & langfristig, grundlegend & Standort, Rechtsform, Geschäftsmodell \\
Laufend & operativ, wiederkehrend & Preis, Marketing, Lagerbestand \\
\bottomrule
\end{tabular}
}
\end{adjustbox}
}



{\small
\begin{adjustbox}{max width=\linewidth}
{\scriptsize
\begin{tabular}{lll}
\toprule
\textbf{Theorie} & \textbf{Frage} & \textbf{Merkmale} \\
\midrule
Normativ & Wie SOLLTE? & rational, modellbasiert (Erwartungswert) \\
Deskriptiv & Wie TATSÄCHLICH? & Biases, begrenzte Rationalität \\
\bottomrule
\end{tabular}
}
\end{adjustbox}
}



\subsection{17. Entscheidungsproblem}

\textbf{Bestandteile}: Alternativen | Umweltzustände | Ergebnisse | Zielgröße

Matrix: Alternative × Umweltzustand → Ergebnis



{\small
\begin{adjustbox}{max width=\linewidth}
{\scriptsize
\begin{tabular}{ll}
\toprule
\textbf{Informationsstand} & \textbf{Bekannt} \\
\midrule
Sicherheit & Konsequenzen vollständig \\
Risiko & Wahrscheinlichkeiten \\
Unsicherheit & weder noch \\
\bottomrule
\end{tabular}
}
\end{adjustbox}
}



\subsection{18. Entscheidung bei Sicherheit}

\textbf{Payment-System}: 200.000 Transaktionen/Jahr; Gesamtkosten = Fixkosten + Gebühr × Transaktionen.



{\small
\begin{adjustbox}{max width=\linewidth}
{\scriptsize
\begin{tabular}{lll}
\toprule
\textbf{Anbieter} & \textbf{Gebühr} & \textbf{Fixkosten} \\
\midrule
A (Basic) & 0,50 € & 5.000 € \\
B (Pro) & 0,35 € & 20.000 € \\
C (Enterprise) & 0,25 € & 50.000 € \\
\bottomrule
\end{tabular}
}
\end{adjustbox}
}



\subsection{19. Entscheidung unter Risiko}

\textbf{Erwartungswert} = Σ (Ergebnis × Wahrscheinlichkeit)



{\small
\begin{adjustbox}{max width=\linewidth}
{\scriptsize
\begin{tabular}{lll}
\toprule
\textbf{Alt.} & \textbf{Erfolg p=0,6} & \textbf{Misserfolg p=0,4} \\
\midrule
A & 100 & 0 \\
B & 200 & −50 \\
\bottomrule
\end{tabular}
}
\end{adjustbox}
}



\subsection{20. Entscheidung unter Unsicherheit}


{\small
\begin{adjustbox}{max width=\linewidth}
{\scriptsize
\begin{tabular}{lll}
\toprule
\textbf{Regel} & \textbf{Logik} & \textbf{Risikohaltung} \\
\midrule
Maximax & bestes Ergebnis je Alt. → Max & risikofreudig \\
Maximin / Wald & schlechtestes je Alt. → Max & risikoavers \\
Minimum-Regret / Savage-Niehans & max. Bedauern je Alt. → Min & Vermeidung Fehlentscheidung \\
Laplace & alle Zustände gleich wahrscheinlich → Ø & neutral \\
Hurwicz & Gewichtung Best/Schlechtest mit λ & flexibel \\
\bottomrule
\end{tabular}
}
\end{adjustbox}
}



\subsection{21. Digitale Entscheidungen}


{\small
\begin{adjustbox}{max width=\linewidth}
{\scriptsize
\begin{tabular}{ll}
\toprule
\textbf{Klassisch} & \textbf{Digital} \\
\midrule
wenig Daten & große Datenmengen \\
statische Modelle & kontinuierliche Optimierung \\
\bottomrule
\end{tabular}
}
\end{adjustbox}
}


\textbf{Anwendungen}: A/B-Testing | Recommendation Systems | Dynamic Pricing


\begin{quote}
Merke: Normative Modelle + Daten = bessere Entscheidungen.
\end{quote}






% ──────────────────────────────────────────────────────────────

\end{multicols}
\end{document}
